\documentclass[11pt]{article}

\usepackage[a4paper,margin=1in]{geometry}
\usepackage{booktabs}
\usepackage{amsmath}
\usepackage{hyperref}
\usepackage{xcolor}
\usepackage{parskip}

\definecolor{codebg}{RGB}{245,245,245}
\newcommand{\code}[1]{\texttt{\small #1}}

\title{\textbf{Group Assignment Report: Question 1} \\[0.3em]
\large Word Segmentation and POS Tagging}
\author{Group 21: Krishna Jhanwar (12341250), Siddharth Rai (12342050),\\
Slok Tulsyan (12342080), Vasu Garg (12342330), Gaurav Kumar (12340790)}
\date{September 8, 2026}

\begin{document}

\maketitle

\section{Overview}

The objective of this assignment was to implement a joint word segmentation and Part-of-Speech (POS) tagging system from scratch. The system had to process a sentence consisting of a single string of characters without spaces, correctly insert word boundaries (segmentation), and assign the correct grammatical category to each segmented word (POS tagging). The task was performed on both English and a morphologically richer language, Spanish, to compare the effectiveness of the models on languages where word endings carry significant grammatical meaning (like gender and number agreement).

\section{Data}

Two datasets were utilized for this task:

\begin{itemize}
    \item \textbf{English:} The Brown Corpus (via the NLTK library), using the Universal Tagset. The dataset was split into an 80/10/10 ratio for Training, Development, and Testing, respectively.
    \item \textbf{Spanish:} The Universal Dependencies Spanish-GSD corpus (UD Spanish-GSD). The standard \code{train}, \code{dev}, and \code{test} \code{.conllu} splits provided by the repository were used.
\end{itemize}

\section{Data Preprocessing \& Morphology Extraction}

The English dataset was relatively straightforward, mapping directly to the simplified Universal Tagset ($\sim$12 base tags). The Spanish dataset required more complex processing. Using the \code{conllu} parser, we extracted the base Universal POS (UPOS) tags alongside their morphological features from the \code{feats} column. We explicitly extracted \code{Gender} and \code{Number} to form \textbf{morphology-aware tags} (e.g., \code{NOUN-Fem-Plur}). This enriched tagset exponentially increased the number of target classes to 86 unique tags, making the classification task substantially harder but allowing the model to learn and enforce grammatical agreement.

\section{Feature Extraction \& Model Architectures}

To tackle the POS tagging task, two distinct model architectures were implemented and compared:

\begin{enumerate}
    \item \textbf{Generative Model (Trigram HMM):} We built a Hidden Markov Model relying on raw probability distributions. It utilizes unigram emission probabilities $P(\text{word} \mid \text{tag})$ and trigram transition probabilities $P(\text{tag}_i \mid \text{tag}_{i-1}, \text{tag}_{i-2})$. Add-$k$ smoothing was applied to handle unseen words and transitions gracefully.
    \item \textbf{Discriminative Model (MEMM):} We implemented a Maximum Entropy Markov Model using scikit-learn's \code{SGDClassifier} optimized for log-loss. For feature extraction, we defined a rich feature set consisting of the current word, prefixes (length 1--3), suffixes (length 1--3), capitalization flags, hyphenation, and the two preceding predicted tags. For English alone, this generated over 160{,}000 sparse features.
\end{enumerate}

For segmentation, a \code{TrigramLanguageModel} was built directly on the training words, capturing the likelihood of word sequences.

\section{Model Training \& Decoding}

\textbf{Training:} The models were trained on the respective 80\% train splits. To handle the massive sparse matrix generated by the MEMM (over 900{,}000 samples and 160{,}000 features for English), we utilized a stochastic gradient descent approach rather than a standard logistic regression solver, preventing memory exhaustion and significantly speeding up convergence.

\textbf{Decoding (Dynamic Programming):} We implemented a \code{ViterbiSegmenter} and a \code{ViterbiPOSTagger}. Because the Spanish morphology-aware tags created an enormous state space (86 tags squared for trigram history), a standard Viterbi search would result in an exponential explosion of states. To resolve this, we introduced a beam-search style optimization into our dynamic programming array, heavily pruning the search space by only retaining the top 50 states at any given timestep $i$. Furthermore, MEMM predictions were batched across all active states simultaneously rather than iteratively, reducing prediction time by orders of magnitude.

\section{Performance \& Baselines}

Before judging the probabilistic models, two simple baselines were established: a Greedy Longest-Match segmenter, and a Most-Frequent-Tag POS tagger.

\subsection*{English Results (Brown Corpus)}

\textbf{Segmentation:} Baseline (85.25\% F1) vs Viterbi Trigram (\textbf{97.08\% F1})

\begin{table}[h]
\centering
\begin{tabular}{lccc}
\toprule
\textbf{POS Model} & \textbf{Seg-Induced Errors} & \textbf{Genuine Errors} & \textbf{Accuracy} \\
\midrule
Baseline (Most Freq) & 415 & 82 & 77.48\% \\
Trigram HMM & 122 & 75 & \textbf{91.07\%} \\
MEMM & 122 & 142 & 88.04\% \\
\bottomrule
\end{tabular}
\end{table}

\subsection*{Spanish Results (UD Spanish-GSD)}

\textbf{Segmentation:} Baseline (80.41\% F1) vs Viterbi Trigram (\textbf{88.27\% F1})

\begin{table}[h]
\centering
\begin{tabular}{lccc}
\toprule
\textbf{POS Model} & \textbf{Seg-Induced Errors} & \textbf{Genuine Errors} & \textbf{Accuracy} \\
\midrule
Baseline (Most Freq) & 594 & 129 & 71.91\% \\
Trigram HMM & 334 & 154 & \textbf{81.04\%} \\
MEMM & 334 & 191 & 79.60\% \\
\bottomrule
\end{tabular}
\end{table}

\section{Evaluation \& Error Analysis}

A custom evaluation module was written to separate downstream errors. By aligning the predicted characters with the gold-standard characters, we categorized every POS tagging error into two groups:

\begin{enumerate}
    \item \textbf{Segmentation-Induced Errors:} The tag is wrong solely because the word boundary was guessed incorrectly (e.g. splitting ``over the'' into ``overt he'').
    \item \textbf{Genuine Tagging Errors:} The word was segmented correctly, but the tagger assigned the wrong POS tag.
\end{enumerate}

Our analysis revealed a critical finding: \textbf{over 60\% of all POS tagging errors were actually segmentation errors in disguise.} This heavily underscores the fragility of a pipeline approach and justifies the necessity of joint decoding.

Furthermore, our models proved that agreement-aware tagging in Spanish was highly effective. Despite the absolute accuracy dropping compared to English (due to the 86-tag state space), the HMM actively learned Spanish grammatical agreements. For example, on the input \code{mispadrespuedenviajar}, it successfully output \code{('mis', 'DET-Plur')}, \code{('padres', 'NOUN-Masc-Plur')}, \code{('pueden', 'AUX-Plur')}, successfully tracking plurality across the sequence.

\section{Conclusion}

The implementation successfully demonstrated the power of dynamic programming and probabilistic modeling over naive greedy algorithms. The Viterbi models vastly outperformed the baselines in both languages. The HMM surprisingly outperformed the MEMM slightly, likely due to the HMM's tighter mathematical coupling with the sequence constraints versus the MEMM's susceptibility to label bias when heavily pruned. Finally, the stark difference between English and Spanish tagging highlighted the complexity of morphologically rich languages, yet showcased our model's capability to learn and enforce deep grammatical agreements organically.

\end{document}
